%Summary:
%\begin{itemize}
%    \item Completely new SiPM readout system developed
%    \item Incorporating modern FPGA data acquisition and FPGA-internal comparators for analog signal discrimination
%    \item Adapted and developed new calibration procedures necessary for FPGA readout system
%    \item Developed lpGBT Mezzanine to interface DT backend
%    \item Developed FPGA firmware for readout and interfacing the lpGBT
%    \item Sucessful modification of DT backend firmware to interface developed readout system
%    \item Installed and commissioned Phase-2 minicrate with DT chamber at test stand in Aachen
%    \item Framework to analyze DT and scintillator data together
%    \item Implemented own triggerless reconstruction code which follows the CMS DT reconstruction
%    \item Successful correlation of measured DT tracks and scintillator hits in space and time
%\end{itemize}
%
%Outlook:
%\begin{itemize}
%    \item Understand open problems
%    \item Implement coincidence IP core (currently still in development)
%    \item Use system to perform various studies with DT chamber, scintillator or even other detectors connected to the developed readout system
%\end{itemize}
\FloatBarrier
\label{sec:SummaryOutlook}

The goal of the thesis was the integration of the scintillation detector into the readout infrastructure of the \ac{DT} chamber at the cosmic test stand. Several tasks were tied to this goal, which have been addressed in this work:
\newline First, the \ac{DT} chamber at the test stand had to be equipped with readout electronics and commissioned. For this, the newly available \ac{CMS} \ac{DT} Phase-2 readout electronics and complementing backend system was installed at the test stand. For the first time in several years, atmospheric muon tracks were reconstructed with the chamber, using a self-implemented triggerless reconstruction algorithm based on the \ac{CMS} \ac{DT} local reconstruction. Although the implemented reconstruction algorithm relies on several simplifications and assumptions, the obtained distributions of the resulting track timing and other properties approximately show the expected behavior. Few systematic features are visible in the fits of the individual \ac{SL} hit patterns, which should be investigated in more detail in the future. Also, the obtained reconstructed track rate of around \SI{85}{\hertz} is smaller than the cosmic expectation.
\newline The second important task was the development of a new readout system for the scintillation detector, which allows the connection to the \ac{DT} backend electronics, and achieve time synchronization with the \ac{DT} chamber. For this purpose, a completely new readout board was conceptualized and produced which hosts and \ac{FPGA} and analog input electronics. The readout system features the \ac{lpGBT} \ac{ASIC}, which provides the actual data and timing interface to the \ac{DT} backend system. The \ac{FPGA} firmware to operate the boards was developed completely in the course of this thesis, and all required calibration methodologies were implemented and validated. During the data taking with the scintillator, systematic features were observed in the recorded hits. While a comprehensive study of the observed effects could not be performed within the timescale of this thesis, potential reasons for this behavior point towards the analog input electronics of the board, which may cause signal reflections and therefore multiple-counting of events. Future investigations of the circuitry and also the \ac{FPGA} firmware are recommended, should the electronics be used beyond this thesis. After applying a dead time to the data to mitigate the observed effects, a lower than expected scintillator rate of around \SI{0.5}{\hertz} per strip could be observed, with strong variations between the channels.
\newline Finally, it was attempted to match the reconstructed \ac{DT} tracks and the scintillator hits in time. From this analysis, the successful synchronization of both detectors can be concluded. The time difference of the matched \ac{DT} tracks and scintillator hits shows a narrow peak width a width of only \SI{12}{\nano\second}. The synchronization between the detectors is expected to be even better than this value, as other effects contribute significantly to the event-by-event timing variations of both detectors. Using the temporally matched tracks and hits, an independent cross-check was performed of the implemented triggerless \ac{DT} reconstruction strategy. The large majority of matched tracks intersect the position of the scintillation detector in the measured geometry, although the need for more precise detector alignment is visible. This result is still considered a success, as the reconstructed tracks generally behave as expected.
\newline With the presented work in this thesis, the cosmic test stand can be declared fully operational. Even though the implemented triggerless \ac{DT} reconstruction strategy should be further tested and optimized, it is evidently already sufficient to generate atmospheric muon tracks with solely the \ac{DT} chamber. In the future, it can be imagined to use the test stand as a reference detector for muons, than can be used to characterize other detectors. This thesis demonstrates the synchronous readout of the chamber with the scintillator, however it is clear that with the gained expertise in the project it is possible to also integrate other particle detectors into the readout system. In this way, one could use the cosmic test stand as a "test bench" for detector research and development.

